\FloatBarrier
\section{Cómo trabajamos}\label{sec:proceso}

\subsection{Un módulo a la vez}\label{sec:incremental}

Construimos el sistema por partes, en el orden en que se necesitan los datos. Primero la base técnica, que es el proyecto, la base de datos, la sesión y los roles. Después los catálogos de clientes y equipos. Luego los lotes y las inspecciones, después el certificado, que usa todo lo anterior, y al final los reportes, la bitácora, las pruebas y los documentos. La \tref{tab:cronologia} resume las fechas.

\begin{table}[H]
  \centering
  \caption{Las etapas principales del proyecto}\label{tab:cronologia}
  \begin{tabular}{@{}ll@{}}
    \toprule
    Fecha & Qué pasó \\
    \midrule
    abril de 2026 & Entrega 1, el documento de arquitectura del equipo \\ %CHECK docs/proceso-desarrollo-es.md says 2026-04-17, docs/llm/context.md says the due date was 2026-03-01
    21 de octubre & Borrador del diseño técnico y prototipo en Node \\
    22 de octubre & Elegimos las herramientas y creamos el proyecto (PR \#1) \\
    4 y 5 de noviembre & Clientes, roles, bitácora, Equipos, Lotes e Inspecciones \\
    11 de noviembre & Certificados con PDF, usuarios y presentación de avances \\
    12 de noviembre & Reportes, Bitácora, menú móvil, CI y guías de prueba \\
    13 de noviembre & Entrega 3 \\
    15 y 16 de noviembre & Correo real, cambios por la revisión del profesor y documentos \\
    \bottomrule
  \end{tabular}
\end{table}

Cada módulo llegó en un pull request propio. Un pull request, o PR, es una petición para unir un cambio al código principal. Cada PR tenía su requisito, sus pruebas y su vista previa en Vercel. Así cada módulo se podía revisar y probar solo, antes de unirlo con los demás.

\subsection{Quién hizo qué}\label{sec:reparto}

Repartimos el trabajo por módulo completo, de la pantalla a la base de datos. Así nadie dependía de otra persona a la mitad de un módulo. \pend{quién hizo cada módulo, a partir de la tabla Organización del equipo de docs/proceso-desarrollo-es.md, confirmada por el equipo}

\pend{cómo se organizó el equipo, con las fechas de las reuniones y el medio que usamos para hablar}

\subsection{El camino de un cambio en GitHub}\label{sec:github}

El código vive en el repositorio del equipo, \urlrepo{}. La \fref{fig:cambio} muestra el camino de un cambio. Cada cambio nace en una rama con un prefijo que dice su tipo, por ejemplo \tt{feat/} para una función nueva o \tt{fix/} para una corrección. Nadie escribe directo en la rama principal, \tt{main}. El cambio llega a \tt{main} solo por un pull request.

\begin{figure}[H]
  \centering
  \resizebox{\linewidth}{!}{% Path of one change, from a branch to production.
% Redrawn from the Git flow in docs/proceso-desarrollo-es.md and docs/llm/runbook.md.
% Drawn at its natural size (no \resizebox), so the text keeps the body font size.
\begin{tikzpicture}[node distance=1.1cm and 0.6cm]
  \tikzset{paso/.style={bloque, text width=3.0cm, minimum height=1.6cm, font=\small}}
  \node[paso] (rama) {\textbf{Rama}\\[2pt]\footnotesize una copia del código para cada cambio};
  \node[paso, right=of rama] (pr) {\textbf{Pull request}\\[2pt]\footnotesize se pide unir el cambio a \tt{main}};
  \node[paso, right=of pr] (ci) {\textbf{Revisión}\\\textbf{automática}\\[2pt]\footnotesize errores, pruebas y construcción};
  \node[paso, right=of ci] (rev) {\textbf{Revisión del equipo}\\[2pt]\footnotesize se lee el cambio y se prueba};

  \node[paso, below=of rev] (merge) {\textbf{Unión a \tt{main}}\\[2pt]\footnotesize el cambio entra al código principal};
  \node[paso, left=of merge, fill=acento!12] (prod) {\textbf{Producción}\\[2pt]\footnotesize Vercel publica la nueva versión};
  \node[paso, below=of rama, fill=white] (prev) {\textbf{Vista previa}\\[2pt]\footnotesize una versión de prueba por rama};

  \draw[flecha] (rama) -- (pr);
  \draw[flecha] (pr) -- (ci);
  \draw[flecha] (ci) -- (rev);
  \draw[flecha] (rev) -- (merge);
  \draw[flecha] (merge) -- (prod);
  \draw[flecha] (rama) -- node[etiq, right] {cada cambio} (prev);
  \draw[flecha, dashed] (ci.north) -- ++(0,0.5) -| node[etiq, pos=0.25, above] {si algo falla, se corrige en la misma rama} (rama.north);
\end{tikzpicture}
}
  \caption{El camino de un cambio, de la rama a producción}\label{fig:cambio}
  \backline{fig:cambio}
\end{figure}

Cada PR pasa por CI. CI, o integración continua, es un programa de GitHub Actions que revisa cada PR solo. Nuestro CI tiene tres trabajos. El primero revisa los tipos, el estilo del código, las pruebas unitarias y la construcción de la aplicación. El segundo corre las pruebas de componentes contra un PostgreSQL de prueba. El tercero abre la aplicación en un navegador con Playwright \parencite{playwright}. Si algo falla, el cambio se corrige en la misma rama.

Los mensajes de commit siguen la convención Conventional Commits \parencite{conventionalcommits} y citan el requisito que cumplen, como en este ejemplo real.

\begin{codigo}
feat(certificados): send certificate emails with Resend (RF-06, D-13)
\end{codigo}

Con esos números se puede ir de cualquier línea de código al requisito que la pide.

Los números de GitHub son estos. El repositorio del equipo tiene 25 pull requests. De ellos, 22 se unieron a \tt{main}, 2 se cerraron sin unir y 1 sigue abierto. También hay 22 issues, que son tareas o preguntas registradas en GitHub. La rama de este reporte tiene 137 commits, del 21 de octubre al 16 de noviembre de 2026. %CHECK counts from gh pr list, gh issue list and git rev-list --count HEAD on 2026-11-16, at commit b90e28c

GitHub no registra ninguna revisión formal en esos PR. De los 22 PR unidos, 6 los unió una persona distinta de quien los escribió, y 16 los unió su propio autor. \pend{cómo revisaba el equipo los PR, por ejemplo en una llamada o en la vista previa, y si alguien más los probaba antes del merge}

\subsection{Lo que cambió por el profesor}\label{sec:profesor}

El profesor revisó el proyecto varias veces mientras lo construíamos. La \tref{tab:profesor} resume lo que pidió y lo que cambiamos.

\begin{table}[H]
  \centering
  \caption{Notas del profesor y lo que cambiamos}\label{tab:profesor}
  \begin{tabularx}{\linewidth}{@{}p{3.2cm}X@{}}
    \toprule
    Nota & Lo que cambiamos \\
    \midrule
    Revisión del 4 de noviembre & Cada cliente tiene un domicilio fiscal y uno o más destinos de entrega. El RFC es único. Los límites de un cliente pueden ser más estrechos que los internacionales, nunca más amplios (D-22, D-23) \\
    Revisión de requisitos & El certificado usa la última inspección del lote, y se bloquea si está incompleta (D-27) \\ %CHECK date of this review, docs/llm/decisions.md only says D-27 was recorded on 2026-11-10
    Presentación del 11 de noviembre & Gráficas de estadísticas y una sección nueva de calidad de la producción en Reportes \\
    Misma presentación & Documentos de arquitectura, de base de datos y del proceso, ordenados en la carpeta \tt{docs/} \\
    \bottomrule
  \end{tabularx}
\end{table}
